Best-channel selections were pooled over the four methods for each dataset, as summarised in Table~\ref{tab:channel_selection}. On Raja (184 selections), the most frequently selected channels were E9 (43.5\%, 80/184) and E22 (37.0\%, 68/184), followed by E3 (8.2\%) and E23 (7.1\%), with selection confined to frontal channels. On Cao2018 (232 selections), the frontopolar channels FP1 (52.6\%, 122/232) and FP2 (26.3\%, 61/232) dominated, while the remaining frontal channels F7, F8, and F3 were each selected in fewer than 8\% of cases. Both datasets therefore exhibited strong frontal and frontopolar dominance, consistent with blink artefacts being largest over frontal sites, although the two corpora used different electrode montages such that channel labels were not directly comparable across datasets. Within-subject consistency was moderate to high, with the median fraction of a subject's sessions selecting that subject's most common channel being 0.667 on Raja and 1.000 on Cao2018. The two proposed variants selected the same channel in 89.1\% of Raja sessions and 82.8\% of Cao2018 sessions, substantially more often than any other method pair, whereas agreement across all four methods on a single channel occurred in only 21.7\% of Raja sessions and 17.2\% of Cao2018 sessions.
